Let us now calculate the second side of the Riemann--Roch formula. For every @@M0@@, we have:
@@M1@@
by (V (6.2.1)). Consequently:
\sgaStableTarget{sga6.e08.eq.5-4-2}@@M2@@
The second side of the Riemann--Roch formula is therefore:
\sgaStableTarget{sga6.e08.eq.5-4-3}@@M3@@
Since @@M4@@ is the class in @@M5@@ of @@M6@@, @@M7@@ is, for @@M8@@, the class in
@@M9@@ of @@M10@@. We thus have:
\sgaStableTarget{sga6.e08.eq.5-4-4}@@M11@@
Finally recall that the elements @@M12@@ are linearly dependent over @@M13@@ (VI 5.6), and satisfy the equation:
\sgaStableTarget{sga6.e08.eq.5-4-5}@@M14@@
and that @@M15@@ are linearly independent over @@M16@@.

One can then write the series @@M17@@ uniquely as a linear combination, with coefficients in
@@M18@@, of the elements @@M19@@. The image of @@M21@@ under @@M20@@ is then, by (5.4.4) and 3.7, its coefficient of @@M22@@. Taking (5.4.1) into account, the Riemann--Roch theorem for @@M23@@ is therefore equivalent to:
\sgaStableTarget{sga6.e08.eq.rr}@@M24@@
This coefficient is the residue of the differential form
@@M25@@
